\href{https://drive.google.com/drive/folders/1le-4uediDAlxbglAtcbnSlPrLuQiPTE3}{\hypertarget{book:C1}{Core Pure I Textbook (C1)}}

\section{Complex Numbers I} \label{complex_numbers_i}

    \subsection{The Imaginary Constant}
        Using only the \glspl{real_number}, there is no number $x$ that satisfies the equation $x^2+1=0$.
        To satisfy this \gls{polynomial}, you must first define a square root of -1.
        This constant is called the \gls{imaginary_constant} $i$, defined such that $i^2=-1$.
        
        Using the \gls{imaginary_constant}, \glspl{complex_number} can be defined as $a+bi$, where $a,b\in\R$.
        The set of all numbers in the form $a+bi$ is called the set of \glspl{complex_number}, written $\C$.

        The functions $\real(z)$ and $\imag(z)$ (sometimes written as $\Re(z)$ and $\Im(z)$) take the \gls{real_part} and \gls{imaginary_part} of a \gls{complex_number} respectively.
        For example,
        \begin{equation} \label{eq:real_part}
            \text{Re}(a+bi)\equiv\Re(z):=a,\ a,b\in\R
        \end{equation}
        \begin{equation} \label{eq:imaginary_part}
            \text{Im}(a+bi)\equiv\Im(z):=bi,\ a,b\in\R
        \end{equation}

    \subsection{Complex Arithmetic}
        \Glspl{complex_number} can be added and multiplied similarly to how binomials would be.
        Given two \glspl{complex_number} $w:=a+bi$ and $z:=c+di$, where $a,b,c,d\in\R$, the following rules define complex addition and multiplication:
        \begin{equation} \label{eqn:complex_addition}
            w+z=(a+bi)+(c+di)=(a+c)+(b+d)i
        \end{equation}
        \begin{equation} \label{eqn:complex_multiplication}
            w\times z=(a+bi)(c+di)=(ac-bd)+(ad+bc)i
        \end{equation}

    \subsection{Complex Conjugates}
        The \gls{complex_conjugate} of a \gls{complex_number} of $z$, written as $z^*$ is $z$ with a flipped sign before the imaginary part.
        \begin{equation}
            (a+bi)^*:=a-bi,\ a,b\in\R
        \end{equation}
        From this definition, we can see that the sum and product of a \gls{complex_number} and its conjugate are both real.
        Let $z:=a+bi$, where $a,b\in\R$.
        \begin{equation} \label{eqn:sum_of_conjugate}
            z+z^*=2a\in\R
        \end{equation}
        \begin{equation} \label{eqn:product_of_conjugate}
            zz^*=a^2+b^2\in\R
        \end{equation}
        Using the multiplication rule in equation \ref{eqn:product_of_conjugate}, a fraction of \glspl{complex_number} can be `realised' by multiplying by the \gls{complex_conjugate} of the denominator over itself.
        \Gls{realising_the_denominator} means rewriting a fraction of \glspl{complex_number} as a fraction with a real denominator.
        \begin{equation} \label{eqn:realise_denominator}
        \begin{split}
            \frac{a+bi}{c+di}&=\frac{a+bi}{c+di}\times\frac{c-di}{c-di}\\
            &=\frac{(a+bi)(c+di)}{c^2+d^2}\\
            &=\frac{(ac-bd)+(ad+bc)i}{c^2+d^2}
        \end{split}
        \end{equation}
    \subsection{Complex Roots of Polynomials}
        All $n$th degree \glspl{polynomial} (\glspl{polynomial} with a power of $x$ up to $x^n$) have exactly $n$ \glspl{root}, counting \glspl{complex_root} and algebraic multiplicities.
        Therefore, if an equation does not cross $y=0$ on a \gls{cartesian_graph}, it must have other, complex \glspl{root}.
        The simplest example of this is the \gls{quadratic} $x^2+1=0$.
        \[x^2=-1\]
        \[x=\pm i\]
        For \glspl{quadratic}, whenever the discriminant ($\Delta:=b^2-4ac$ for a \gls{quadratic} $ax^2+bx+c$) is less than $0$, the \gls{quadratic} has \glspl{complex_root}, since the quadratic formula (\ref{eq:quadratic_formula}) can be rewritten as
        \[x=\frac{-b\pm\sqrt{\Delta}}{2a}\]
        which would give a \gls{complex_number} whenever $\Delta<0$, since $\sqrt{x}$ is complex iff $x<0$ ($x\in\R$).
        
        If a \gls{polynomial} has a \gls{complex_root} $\alpha$, it must also have the \gls{complex_conjugate} \gls{root} $\alpha^*$.
        For example, if a \gls{quadratic} has a \gls{root} $\alpha=2+3i$, then the other \gls{root} must be $\alpha^*=2-3i$.
        For a \gls{cubic}, there can only be one pair of \glspl{root}, thus the third \gls{root} cannot be a \gls{complex_root}, since it would have to have a fourth \gls{root} to form a \gls{conjugate_pair}.
        Thus, the third \gls{root} of a \gls{cubic} must a \gls{real_number}.

    \subsection{Example Questions}
        \question{q:1.1}{C1}{1A/8}{}{
            Given that $z=a+bi$ and $w=a-bi$, $a,b\in\R$, show that:
            \begin{enumerate}
                \item $z+w$ is always real
                \item $z-w$ is always imaginary
            \end{enumerate}
        }
        \answer{q:1.1}{
            \begin{enumerate}
                \item $z+w$ can be expanded to $(a+bi)+(a-bi)$, which simplifies to $2a+bi-bi=2a$. Since $a\in\R$, $2a\in\R$, thus $z+w\in\R$.\qed
                \item $z-w$ can be expanded to $(a+bi)-(a-bi)$, which simplifies to $a-a+2bi=2bi$. Since $b\in\R$, $2b\in\R$, thus $2bi$ is purely imaginary.\qed
            \end{enumerate}
        }

        \question{q:1.2}{C1}{1C/9}{}{
            \begin{enumerate}
                \item Given that $i^1=i$ and $i^2=-1$, write $i^3$ and $i^4$ in their simplest forms.
                \item Write $i^5$, $i^6$, $i^7$ and $i^8$ in their simplest forms.
                \item Write down the value of:
                \begin{enumerate}
                    \item $i^{100}$
                    \item $i^{253}$
                    \item $i^{301}$
                \end{enumerate}
            \end{enumerate}
        }
        \answer{q:1.2}{
            \begin{enumerate}
                \item Since $i^2=-1$, $i^3$ must be $i^2\times i=-1\times i=\boxed{-i}$. Then, $i^4$ is $i^3\times i=-i\times i=-(i^2)=-(-1)=\boxed{1}$.
                \item $i^5$ is $i^4\times i=1\times i=\boxed{i}$; $i^6$ is $i^5\times i=\boxed{-1}$; $i^7=i^6\times i=\boxed{-i}$; $i^8=i^7\times i=\boxed{1}$.
                \item From the pattern in parts a and b, it is clear that $i^{n+4k}=i^n$ for all $k\in\Z$, thus $i^n=i^{n\mod 4}$.
                \begin{enumerate}
                    \item $100\mod4=0$, so $i^{100}=i^0=\boxed{1}$.
                    \item $253\mod4=1$, so $i^{253}=i^1=\boxed{i}$.
                    \item $301\mod4=1$, so $i^{301}=i^1=\boxed{i}$.
                \end{enumerate}
            \end{enumerate}
        }

        \question{q:1.3}{C1}{1C/Challenge}{}{
            \begin{enumerate}
                \item Expand $(a+bi)^2$
                \item Hence, or otherwise, find $\sqrt{40-42i}$, giving your answer in the form $a-bi$, where $a$ and $b$ are positive integers.
            \end{enumerate}
        }
        \answer{q:1.3}{
            \begin{enumerate}
                \item $(a+bi)^2=\boxed{a^2+2abi-b^2}$.
                \item Let $z=a-bi=\sqrt{40-42i}$.
                Thus, $z^2=40-42i=a^2-2abi-b^2$.
                This shows that $40=a^2-b^2$ and $-42=-2ab$, since the real and imaginary terms must match up.
                Rearranging $-42=2ab$ in terms of $a$ gives:
                \begin{equation*}
                    \begin{aligned}
                        -42&=-2ab \\
                        21&=ab \\
                        a&=\frac{21}{b}
                    \end{aligned}
                \end{equation*}
                Which can then be substituted into $40=a^2+b^2$:
                \begin{equation*}
                    \begin{aligned}
                        40&=(\frac{21}{b})^2-b^2 \\
                        40&=\frac{441}{b^2}-b^2 \\
                        40b^2&=441-b^4
                    \end{aligned}
                \end{equation*}
                This is now a quadratic in terms of $b^2$, so we'll define $\beta=b^2$, and it can now be solved using the quadratic equation:
                \begin{equation*}
                    \begin{aligned}
                        b^4+40b^2-441&=0 \\
                        \beta^2+40\beta-441&=0 \\
                        \beta&=\frac{-40\pm\sqrt{40^2-4(-441)}}{2} \\
                        &=\frac{-40\pm\sqrt{3364}}{2} \\
                        &=\frac{-40\pm58}{2} \\
                        &=-20\pm29 \\
                        \beta=9\;\text{or}\;\beta=-49 \\
                        b=\pm\sqrt{\beta} \\
                    \end{aligned}
                \end{equation*}
                Thus, $b=3$ or $b=-3$ or $b=7i$ or $b=-7i$.
                Since the question stipulates that \textit{``$a$ and $b$ are positive integers''}, $b$ must be 3.
                Therefore, $a=\frac{21}{3}=7$.
                So $\sqrt{40-42i}=\boxed{7-3i}$.
            \end{enumerate}
        }

    \newpage

\section{Argand Diagrams} \label{argand_diagrams}
    \subsection{Complex Numbers on Argand Diagrams}
        An \gls{argand_diagram} can be thought of similarly to a \gls{cartesian_graph}; \glspl{complex_number} in the form $x+iy$ can be plotted like $(x, y)$ can on a \gls{cartesian_graph}.
        Instead of an $x$ axis and a $y$ axis, there are instead a real axis and an imaginary axis, used to plot the \gls{real_part} and \gls{imaginary_part} of \glspl{complex_number}.
        An \gls{argand_diagram} looks like this:

        \begin{center}
        \begin{tikzpicture}[scale=1.5, >=stealth] % From Gemini
            % Draw axes
            \draw[->] (-0.5, 0) -- (4, 0) node[right] {$\real(z)$};
            \draw[->] (0, -0.5) -- (0, 3) node[above] {$\imag(z)$};
            
            % Define coordinates for complex number z = a + bi (e.g., z = 3 + 2i)
            \coordinate (O) at (0,0);
            \coordinate (Z) at (3,2);
            \coordinate (X) at (3,0);
            
            % Draw vector from origin to z
            %\draw[thick, red] (O) -- (Z) node[midway, above left] {$z = a+bi$};
            \draw[thick, red] (O) -- (Z);
            \draw[red] (Z) node[above] {$z=a+bi$};
            % Projection lines to axes
            \draw[dashed, gray] (Z) -- (X) node[below] {$a$};
            \draw[dashed, gray] (Z) -- (0,2) node[left] {$b$};
            
            % Plot point at z
            \filldraw[red] (Z) circle (1.5pt);
            
            % Draw angle theta
            %\draw (0.5,0) arc[start angle=0, end angle=33.7, radius=0.5];
            %\node at (0.7, 0.25) {$\theta$};
        \end{tikzpicture}
        \end{center}

        Where the point $z=a+bi$ is drawn where the point $(a, b)$ would be drawn on a \gls{cartesian_graph}.
        
    \subsection{Modulus-Argument Form}
        The \gls{modulus_argument_form} of a \gls{complex_number} is similar to the polar form of a point on a \gls{cartesian_graph}.
        It defines the point using its distance from the origin, and the angle it makes with the positive half of the real axis, in \glspl{radian}.
        The \gls{modulus} of $z$ is written as $|z|$, which is calculated using Pythagoras' Theorem:
        \begin{equation}
            |z|=\sqrt{\real^2(z)+\imag^2(z)}
        \end{equation}
        where $\real^2(z)$ and $\imag^2(z)$ mean $(\real(z))^2$ and $(\imag(z))^2$ respectively.
        The angle that $z$ makes with the positive real number line is called the \gls{argument}, and is positive, going anti-clockwise if $\imag(z)\ge0$, and negative, going clockwise if $\imag(z)<0$.
        The \gls{argument} of $z$ is written as $\arg(z)$.

        An \gls{argand_diagram} is split into four \glspl{quadrant}, called \gls{quadrant} 1, 2, 3 and 4:

        \begin{center}
        \begin{tikzpicture}[scale=1.5, >=stealth]
            \draw[->] (-2, 0) -- (2, 0) node[right] {$\real(z)$};
            \draw[->] (0, -2) -- (0, 2) node[above] {$\imag(z)$};
            
            \coordinate (Q1) at (1,1);
            \coordinate (Q2) at (-1,1);
            \coordinate (Q3) at (-1,-1);
            \coordinate (Q4) at (1,-1);
            
            \draw (Q1) node {Quadrant 1};
            \draw (Q2) node {Quadrant 2};
            \draw (Q3) node {Quadrant 3};
            \draw (Q4) node {Quadrant 4};
        \end{tikzpicture}
        \end{center}


        For values of $z$ in the first and fourth \glspl{quadrant} ($\real(z)>0$):
        \begin{equation}
            \arg(z)=\arctan\left(\frac{\imag(z)}{\real(z)}\right)
        \end{equation}
        For values of $z$ in the second \gls{quadrant}:
        \begin{equation}
            \arg(z)=\arctan\left(\frac{\imag(z)}{\real(z)}\right)+\pi
        \end{equation}
        For values of $z$ in the third \gls{quadrant}:
        \begin{equation}
            \arg(z)=\arctan\left(\frac{\imag(z)}{\real(z)}\right)-\pi
        \end{equation}

        On an \gls{argand_diagram}, the \gls{complex_conjugate} of a number is itself mirror in the line $\imag(z)=0$ (see subsection \ref{ssec:loci_and_regions_in_argand_diagrams}):

        \begin{center}
        \begin{tikzpicture}[scale=1.5, >=stealth]
            \draw[->] (-4, 0) -- (4, 0) node[right] {$\real(z)$};
            \draw[->] (0, -3) -- (0, 3) node[above] {$\imag(z)$};
            
            \coordinate (O) at (0,0);
            \coordinate (Z) at (3,2);
            \coordinate (conj_Z) at (3,-2);
            \coordinate (neg_Z) at (-3,-2);
            \coordinate (neg_conj_Z) at (-3,2);
            
            \draw[thick, red] (O) -- (Z);
            \draw[thick, red] (O) -- (conj_Z);
            \draw[thick, red] (O) -- (neg_Z);
            \draw[thick, red] (O) -- (neg_conj_Z);

            \draw[red] (Z) node[above] {$z=a+bi$};
            \draw[red] (conj_Z) node[below] {$z^*=a-bi$};
            \draw[red] (neg_Z) node[below] {$-z=-a-bi$};
            \draw[red] (neg_conj_Z) node[above] {$-z^*=-a+bi$};

            \filldraw[red] (Z) circle (1.5pt);
            \filldraw[red] (conj_Z) circle (1.5pt);
            \filldraw[red] (neg_Z) circle (1.5pt);
            \filldraw[red] (neg_conj_Z) circle (1.5pt);
        \end{tikzpicture}
        \end{center}

    \subsection{Loci and Regions in Argand Diagrams} \label{ssec:loci_and_regions_in_argand_diagrams}

    \subsection{Example Questions}

    \newpage

\section{Series I} \label{series_i}

\section{Roots of Polynomials} \label{roots_of_polynomials}
    \subsection{Sums of Roots}
        The \glspl{root} of \glspl{polynomial} can be related using Vieta's formulae (see section \ref{vietas_formulae}).
        The \glspl{root} of a general \gls{polynomial} are written as $\alpha$, $\beta$, $\gamma$, $\delta$, etc., and the coefficients are written as $a$, $b$, $c$, $d$, $e$, etc., where $a$ is the coefficient of the greatest power of the \gls{polynomial}, $b$ is the coefficient of the next greatest power, etc.
        The sum of all \glspl{root} is then written as $\sum\alpha\equiv\alpha+\beta+\cdots$.
        Multiple \glspl{root} can be multiplied together and then summed, for example the sum of all pairs of \glspl{root} ($\alpha\beta+\alpha\gamma+\cdots$) would be written as $\sum\alpha\beta$.
        This would be the sum of \textit{all} possible pairs of \glspl{root}.

    \subsection{Linear Transformations of Roots}

\newpage
\section{Volumes of Revolution} \label{volumes_of_revolution}
\newpage
\section{Matrices} \label{matricies}
\newpage
\section{Linear Transformations} \label{linear_transformations}
\newpage
\section{Proof by Induction} \label{proof_by_induction}
\newpage
\section{Vectors} \label{core_vectors}
\newpage